\documentclass[11pt,letterpaper]{article}
\usepackage[margin=1in]{geometry}
\usepackage{graphicx}
\usepackage{hyperref}
\usepackage{booktabs}
\usepackage{xcolor}
\usepackage{amsmath}
\usepackage{amssymb}
\usepackage{tikz}
\usetikzlibrary{arrows.meta, positioning, calc, fit}

\definecolor{codebg}{HTML}{F5F5F5}
\definecolor{hot}{HTML}{EF4444}
\definecolor{warm}{HTML}{F59E0B}
\definecolor{cold}{HTML}{3B82F6}
\definecolor{gapred}{HTML}{DC2626}

\title{\textbf{Adaptive RDMA Optimization for Distributed LLM Workloads}\\[0.5em]
\large Update 11: Mechanistic Interpretability Meets Distributed Tensor Transport\\[0.3em]
\normalsize Insights from the Transformer Circuits Thread}
\author{Nathan Gopee\\
MS Computer Science, SUNY New Paltz\\
Advisors: Dr.\ Wafi Danesh \& Dr.\ Ashley Suchy}
\date{February 2026}

\begin{document}
\maketitle

\begin{abstract}
\noindent
This update bridges Anthropic's Transformer Circuits research program (2021--2025) with the adaptive RDMA middleware developed in Updates~2--10. By analyzing how transformers actually compute internally---which layers matter, how features compose across layers, and how sparse activations really are---we identify concrete gaps in the current middleware design and propose circuit-aware optimizations. The key finding is that mechanistic interpretability provides the missing ``why'' behind our WFA tensor classifier: feature sparsity, cross-layer composition, and induction head access patterns give principled reasons for the hot/warm/cold routing decisions that Update~10 makes heuristically.
\end{abstract}

\tableofcontents
\newpage

%=============================================================================
\section{The Transformer Circuits Research Program}
\label{sec:circuits-overview}
%=============================================================================

Anthropic's Transformer Circuits Thread is a multi-year research effort to reverse-engineer the computational mechanisms inside transformer language models. Starting with a mathematical framework for small attention-only models~\cite{elhage2021framework} and progressing through dictionary learning on production-scale systems~\cite{templeton2024scaling}, the program now includes circuit tracing on Claude 3.5 Haiku~\cite{lindsey2025biology} with 30 million learned features.

The research is organized around a core observation: individual neurons in transformers are \textbf{polysemantic}---a single neuron may respond to unrelated concepts (DNA sequences, legal language, Korean text). This is not a bug but an efficient encoding strategy called \textbf{superposition}~\cite{elhage2022superposition}, where models store more concepts than they have dimensions by exploiting the near-orthogonality of sparse feature directions.

\subsection{Why This Matters for RDMA}

The Circuits research answers questions that directly affect tensor transport decisions:

\begin{enumerate}
    \item \textbf{Which layers matter for a given computation?} Circuit tracing reveals that features routinely skip layers---a concept activated at layer 3 may directly influence layer 12 with no meaningful computation in between~\cite{ameisen2025tracing}.
    \item \textbf{How sparse are activations really?} On any given token, only $\sim$235 features out of 30 million are active (L0 = 235 in Claude 3.5 Haiku)~\cite{lindsey2025biology}. At the SAE level, $<$300 features per token out of millions~\cite{templeton2024scaling}.
    \item \textbf{What access patterns do attention heads create?} Induction heads generate content-dependent, long-range lookback into the KV cache~\cite{olsson2022induction}, creating the non-uniform access patterns our WFA classifier observes.
    \item \textbf{How do precision requirements vary across computation?} Parallel internal pathways have different precision needs---approximate estimation streams tolerate quantization while precise digit-calculation streams do not~\cite{lindsey2025biology}.
\end{enumerate}

%=============================================================================
\section{Key Findings from the Circuits Thread}
\label{sec:findings}
%=============================================================================

\subsection{The Residual Stream as a Communication Bus}

The mathematical framework~\cite{elhage2021framework} reconceptualizes the transformer not as a stack of layers but as a \textbf{residual stream}: a high-dimensional vector per token position that every component reads from and writes to additively. Each attention head and MLP layer contributes its output to a running sum:
\begin{equation}
    \mathbf{x}_{\text{final}} = \mathbf{W}_E + \mathbf{W}_{\text{pos}} + \sum_{\ell} \sum_h \text{head}_{\ell,h}(\mathbf{x}) + \sum_{\ell} \text{MLP}_\ell(\mathbf{x})
\end{equation}

Because every component's contribution is additive, the output logits decompose into a linear sum of each component's individual effect. This enables \textbf{direct logit attribution}: measuring how much each head or MLP contributes to the final prediction for any input.

The residual stream has \textbf{no privileged basis}~\cite{anthropic2023privileged}---individual dimensions are meaningless; only directions carry semantic content. However, Adam optimizer normalization creates stable outlier dimensions with 20$\times$ magnitude, which has direct implications for quantization (Section~\ref{sec:gaps}).

\subsection{QK and OV Circuits: Two Independent Computations per Head}

Each attention head implements two mathematically independent circuits~\cite{elhage2021framework}:

\begin{itemize}
    \item \textbf{QK circuit} ($\mathbf{W}_Q \mathbf{W}_K^\top$): Determines \emph{where to attend}. A rank-$d_{\text{head}}$ bilinear form over token pairs.
    \item \textbf{OV circuit} ($\mathbf{W}_V \mathbf{W}_O$): Determines \emph{what to move}. Maps attended-to information into the residual stream.
\end{itemize}

The full token-level QK circuit $\mathbf{W}_E^\top \mathbf{W}_{QK} \mathbf{W}_E$ can be precomputed to predict attention patterns from token identity alone---a signal for KV cache prefetching.

\subsection{Composition: Virtual Attention Heads}

Multi-layer transformers gain power through \textbf{head composition}~\cite{elhage2021framework}. A layer-0 head's output enters the residual stream and modifies the queries, keys, or values of layer-1 heads, creating three composition types:

\begin{table}[htbp]
\centering
\caption{Attention head composition types and their implications for RDMA transfer scheduling}
\label{tab:composition}
\begin{tabular}{@{}llll@{}}
\toprule
\textbf{Type} & \textbf{Mechanism} & \textbf{Effective Result} & \textbf{RDMA Implication} \\
\midrule
Q-composition & Head 0 output $\to$ Head 1 query & Dynamic search criteria & Transfer must precede QK compute \\
K-composition & Head 0 output $\to$ Head 1 key & Enriched source tokens & KV cache depends on prior layer \\
V-composition & Head 0 output $\to$ Head 1 value & Chained information move & Composition chain = latency-critical \\
\bottomrule
\end{tabular}
\end{table}

Composition strength is measurable: $\|\mathbf{W}_{OV}^{h_0} \cdot \mathbf{W}_K^{h_1}\|_F$ quantifies how strongly head $h_0$ influences head $h_1$'s attention pattern. This gives a \textbf{computable, static metric} for RDMA transfer priority between pipeline stages.

\subsection{Induction Heads and In-Context Learning}

Induction heads~\cite{olsson2022induction} implement pattern completion: given $[A][B]\ldots[A]$, predict $[B]$. They emerge through a sharp phase transition during training and are the primary mechanism for in-context learning.

The circuit requires exactly two heads across two layers:
\begin{enumerate}
    \item \textbf{Previous-token head} (layer 0): Writes predecessor identity into the residual stream via positional attention.
    \item \textbf{Induction head} (layer 1): Uses K-composition to match current token against predecessor annotations, then copies the following token.
\end{enumerate}

Induction heads create \textbf{content-dependent, long-range, sparse lookback} into the KV cache---exactly the non-local access pattern that our WFA classifier must handle. Only 1--3\% of heads are induction/retrieval heads, but they require full KV cache access. This asymmetry validates the hot/warm/cold classification: induction head KV entries are HOT (dedicated QP, polled completions), while non-retrieval heads can use compressed cache (COLD, batched).

\subsection{Superposition and Extreme Sparsity}

The superposition hypothesis~\cite{elhage2022superposition} is proven in toy models: networks store $n \gg m$ features in $m$ dimensions using near-orthogonal directions. This works because features are sparse---interference is rarely realized when few features are simultaneously active.

\begin{table}[htbp]
\centering
\caption{Activation sparsity across Transformer Circuits publications}
\label{tab:sparsity}
\begin{tabular}{@{}llrl@{}}
\toprule
\textbf{Model} & \textbf{Method} & \textbf{Active / Total} & \textbf{Sparsity} \\
\midrule
1-layer transformer~\cite{bricken2023monosemantic} & SAE (8$\times$) & 30 / 4,096 & 99.3\% \\
Claude 3 Sonnet~\cite{templeton2024scaling} & SAE (34M) & $<$300 / 34M & 99.999\% \\
Claude 3.5 Haiku~\cite{lindsey2025biology} & CLT (30M) & 235 / 30M & 99.999\% \\
\bottomrule
\end{tabular}
\end{table}

The phase diagram from~\cite{elhage2022superposition} identifies sharp transitions between three representation states:
\begin{itemize}
    \item \textbf{Dedicated}: Feature gets its own dimension. Full precision needed.
    \item \textbf{Superposition}: Feature shares dimensions. Tolerates some quantization noise (already experiencing interference).
    \item \textbf{Absent}: Feature not represented. No transfer needed.
\end{itemize}

The transitions are discontinuous (first-order), meaning compression quality does not degrade smoothly---there are natural breakpoints where increasing compression has minimal impact until a phase boundary is crossed.

\subsection{Cross-Layer Feature Persistence}

Crosscoders~\cite{anthropic2024crosscoders} prove that significant redundant structure exists across layers: features computed at one layer persist in the residual stream through many subsequent layers. A single crosscoder feature spanning 10 layers replaces 10 redundant per-layer SAE features.

This directly supports \textbf{delta-based tensor transport}: send a base representation from the first layer where a feature activates, then transmit only per-layer deltas for subsequent layers. The cross-layer redundancy means most deltas are near-zero.

\subsection{Feature Steering and Causal Significance}

Scaling Monosemanticity~\cite{templeton2024scaling} demonstrates that clamping individual SAE features predictably modifies model behavior---from making Claude obsess about the Golden Gate Bridge to overriding RLHF safety training. Features are not correlational artifacts; they are causally significant computational units.

For our SAE feature steering work (Update~10 Section~12), this validates that the sparse feature representation we use for gradient-free behavioral modification is operating on the same fundamental units that Anthropic's mechanistic interpretability identifies.

\subsection{The Biology of Computation}

Circuit tracing on Claude 3.5 Haiku~\cite{lindsey2025biology} reveals how a production model actually computes. Key architectural patterns:

\begin{itemize}
    \item \textbf{Multi-hop reasoning}: Intermediate representations (``Dallas'' $\to$ ``Texas'' $\to$ ``Austin'') traverse the residual stream across layers. Pipeline-parallel partitioning must preserve these chains.
    \item \textbf{Planning features}: The model holds candidate outputs before generating (rhyme targets, diagnostic hypotheses). Early layers predict late-layer needs---a signal for RDMA prefetching.
    \item \textbf{Parallel precision pathways}: Arithmetic uses simultaneous approximate-estimation and precise-digit streams. Different pathways in the \emph{same layer} have different precision requirements.
    \item \textbf{Default refusal with override}: Safety operates as an always-on refusal circuit suppressed by ``known entity'' features. Hallucination occurs when the override fires incorrectly.
\end{itemize}

%=============================================================================
\section{Gap Analysis: What the Middleware is Missing}
\label{sec:gaps}
%=============================================================================

Mapping Circuits findings against the current Update~10 middleware design reveals five concrete gaps.

\subsection{Gap 1: WFA Classification is Heuristic, Not Circuit-Aware}

\textbf{Current state}: The WFA classifier uses access-frequency thresholds ($\theta_1 = 5$, $\theta_2 = 20$) and idle timeouts ($t_{\text{idle}} = 0.5$s) to classify tensors as hot/warm/cold. These thresholds are hand-tuned.

\textbf{What Circuits reveals}: Tensor importance is not just about access frequency. A tensor accessed 5 times might be on a critical composition chain (K-composition feeding an induction head) while a tensor accessed 50 times might be a redundant copy of a cross-layer feature.

\textbf{Gap}: The WFA lacks a \textbf{composition-awareness} signal. Two metrics from the framework could augment it:

\begin{equation}
    \text{priority}(h) = \underbrace{\|\mathbf{W}_E \cdot \mathbf{W}_{OV}^h \cdot \mathbf{W}_U\|_\sigma}_{\text{direct importance (spectral norm)}} + \sum_{h' \in \text{downstream}} \underbrace{\|\mathbf{W}_{OV}^h \cdot \mathbf{W}_K^{h'}\|_F}_{\text{composition strength}}
\end{equation}

Heads with high composition strength but low direct attribution are ``invisible but critical''---they would be classified WARM or COLD by access frequency alone but should be HOT.

\subsection{Gap 2: Precision Management Ignores Superposition Phase}

\textbf{Current state}: The precision pipeline (Update~10 Section~10) adapts between FP16/BF16/FP8/MXFP4 based on tensor type (weights, activations, gradients) and hardware (V100 FP16, 5070~Ti BF16).

\textbf{What Circuits reveals}: Precision requirements depend on the \textbf{superposition phase} of the features being transferred:
\begin{itemize}
    \item \textbf{Dedicated features}: Zero interference margin. Quantization noise directly corrupts the signal. Needs FP16+.
    \item \textbf{Superposition features}: Already experiencing interference from overlapping directions. Quantization noise is a small additional cost. Tolerates FP8 or MXFP4.
    \item \textbf{Dense important features} (attention, residual stream): Many simultaneous features means low interference tolerance. High precision required.
    \item \textbf{Sparse MLP activations} (post-ReLU): 50--90\% zeros. Sparsity encoding + aggressive quantization on nonzeros.
\end{itemize}

Additionally, privileged basis outlier dimensions~\cite{anthropic2023privileged} with 20$\times$ magnitude devastate uniform quantization. Channel-wise or per-dimension scaling is needed.

\textbf{Gap}: The precision pipeline treats all activations at a given layer uniformly. It should differentiate by \textbf{feature phase} and \textbf{dimensional outliers}.

\subsection{Gap 3: No Sparse Activation Compression}

\textbf{Current state}: Activation tensors are transferred as dense buffers over RDMA, optionally with MXFP4 quantization (4$\times$ compression).

\textbf{What Circuits reveals}: In the SAE feature basis, activations are 99.3--99.999\% sparse. Even without SAE decomposition, post-ReLU MLP activations are 50--90\% zero. Transferring dense tensors wastes bandwidth on zeros.

\textbf{Gap}: The middleware has no \textbf{sparsity-aware transfer mode}. Options:
\begin{itemize}
    \item \textbf{Zero-skip encoding}: Transmit only nonzero indices + values. For 90\% sparse tensors, this is a 10$\times$ compression before any quantization.
    \item \textbf{Feature-basis projection}: Project activations into a sparse basis before transfer, transmit $\sim$300 (index, value) pairs instead of 12K+ dense floats ($\sim$20$\times$ compression at FP16).
    \item \textbf{Combined}: Sparse encoding + MXFP4 on nonzero values = 40--80$\times$ theoretical compression.
\end{itemize}

\subsection{Gap 4: KV Cache Prefetching is Pattern-Blind}

\textbf{Current state}: The prefetch FSM (Update~10 Section~10.5) uses a simple state machine (Idle $\to$ Prefetching $\to$ Ready) driven by the WFA classifier. It does not know \emph{which} KV cache entries will be needed.

\textbf{What Circuits reveals}: KV cache access patterns are predictable from circuit structure:
\begin{itemize}
    \item \textbf{Previous-token heads}: Always access position $t-1$. Stride-1 prefetch.
    \item \textbf{Induction heads}: Access positions where $\text{token}_{j-1} = \text{current token}$. Content-indexed prefetch.
    \item \textbf{BOS/global heads}: Always access position 0. Pin in cache.
\end{itemize}

Furthermore, only 1--3\% of heads are retrieval heads requiring full cache access. The rest can use compressed or partial KV cache~\cite{olsson2022induction}.

\textbf{Gap}: The prefetch module should maintain a \textbf{head-type registry} (previous-token, induction, global, local-window) and issue targeted RDMA reads based on head type rather than blanket prefetching.

\subsection{Gap 5: No Cross-Layer Delta Compression}

\textbf{Current state}: Each pipeline stage transfers full activation tensors to the next stage.

\textbf{What Circuits reveals}: Crosscoders~\cite{anthropic2024crosscoders} show that features persist across many layers with minimal change. The residual stream at layer $\ell+1$ is layer $\ell$'s stream plus a small delta from the MLP and attention heads.

\textbf{Gap}: For pipeline-parallel inference where consecutive layers are on different nodes, sending only the \textbf{delta} (the head and MLP output) rather than the full residual stream could significantly reduce transfer size. The delta is bounded by:
\begin{equation}
    \|\delta_\ell\| = \|\text{head}_\ell(\mathbf{x}) + \text{MLP}_\ell(\mathbf{x})\| \ll \|\mathbf{x}_\ell\|
\end{equation}
particularly in later layers where the residual stream norm grows while per-layer contributions remain bounded.

%=============================================================================
\section{Proposed Circuit-Aware Enhancements}
\label{sec:proposals}
%=============================================================================

\subsection{Enhancement 1: Composition-Weighted WFA}

Augment the existing WFA state machine with a static composition weight per tensor, computed once from model weights:

\begin{equation}
    \text{WFA\_class}(t) = f\big(\text{access\_count}(t),\; \text{elapsed}(t),\; \underbrace{w_{\text{comp}}(t)}_{\text{new: static weight}}\big)
\end{equation}

The composition weight biases classification upward for tensors on critical composition chains. A tensor with low access count but high $w_{\text{comp}}$ stays classified as WARM or HOT.

\subsection{Enhancement 2: Phase-Aware Precision Allocation}

Extend the precision pipeline with three new signals:

\begin{enumerate}
    \item \textbf{Sparsity ratio}: Measured at runtime per tensor. High sparsity $\to$ allow MXFP4.
    \item \textbf{Outlier magnitude}: Per-dimension max/mean ratio. High outliers $\to$ per-channel scaling before quantization.
    \item \textbf{Composition depth}: Number of downstream composition hops. Deep chains $\to$ higher precision (quantization errors compound through $\mathbf{W}_{OV}^{h_0} \cdot \mathbf{W}_{OV}^{h_1}$).
\end{enumerate}

\subsection{Enhancement 3: Sparse RDMA Transfer Mode}

Add a sparse transfer path to the transport layer:

\begin{enumerate}
    \item Sender scans activation tensor, extracts nonzero indices and values.
    \item Packs into a sparse message: \texttt{[count:u32][indices:u16[]][values:fp16[]]}.
    \item Receiver reconstructs dense tensor from sparse message.
\end{enumerate}

For a $d = 12{,}288$ activation vector at 90\% sparsity: dense FP16 = 24~KB, sparse = $\sim$4.8~KB (5$\times$ compression). At 99\% sparsity (SAE basis): sparse = $\sim$0.5~KB (48$\times$).

\subsection{Enhancement 4: Head-Type-Aware Prefetching}

Classify each attention head by type (from weight analysis, computed once):

\begin{table}[htbp]
\centering
\caption{Head-type-aware KV cache prefetch strategy}
\label{tab:prefetch}
\begin{tabular}{@{}llll@{}}
\toprule
\textbf{Head Type} & \textbf{Detection Method} & \textbf{Prefetch Strategy} & \textbf{Cache Policy} \\
\midrule
Previous-token & Positional QK diagonal & Stride-1 sequential & Sliding window \\
Induction & High K-composition with prev-tok head & Content-indexed lookup & Full retention \\
BOS/global & Uniform attention to pos 0 & Pin position 0 & Never evict \\
Local window & Banded QK structure & Fixed-window prefetch & LRU within window \\
Semantic & Content-based QK & Reactive (no prefetch) & Frequency-based evict \\
\bottomrule
\end{tabular}
\end{table}

\subsection{Enhancement 5: Delta-Mode Pipeline Transfer}

For pipeline-parallel inference, the receiving node maintains the previous residual stream state. The sender transmits only the per-layer delta:

\begin{equation}
    \text{transfer}(\ell) = \text{head}_\ell(\mathbf{x}) + \text{MLP}_\ell(\mathbf{x})
\end{equation}

The receiver reconstructs: $\mathbf{x}_{\ell+1} = \mathbf{x}_\ell + \text{transfer}(\ell)$. Combined with sparse encoding on the delta (which is sparser than the full residual stream), this compounds bandwidth savings.

%=============================================================================
\section{Mapping Circuits Findings to Existing Middleware}
\label{sec:mapping}
%=============================================================================

\begin{table}[htbp]
\centering
\caption{Transformer Circuits findings mapped to \texttt{tensor\_rdma} modules}
\label{tab:mapping}
\footnotesize
\begin{tabular}{@{}p{3.8cm}p{3cm}p{2.5cm}p{3.5cm}@{}}
\toprule
\textbf{Circuits Finding} & \textbf{Source} & \textbf{Module} & \textbf{Application} \\
\midrule
Extreme activation sparsity (L0=235/30M) & Biology~\cite{lindsey2025biology} & \texttt{transport.rs} & Sparse transfer encoding \\
Composition strength measurable from weights & Framework~\cite{elhage2021framework} & \texttt{config.rs} & Static priority computation \\
Induction heads = content-indexed KV access & Induction~\cite{olsson2022induction} & \texttt{prefetch.rs} & Content-aware prefetch \\
Phase transitions in representation & Superposition~\cite{elhage2022superposition} & \texttt{precision.rs} & Phase-aware quantization \\
Cross-layer feature persistence & Crosscoders~\cite{anthropic2024crosscoders} & \texttt{rdma\_cache.rs} & Delta compression, dedup \\
Parallel precision pathways & Biology~\cite{lindsey2025biology} & \texttt{precision.rs} & Per-pathway precision \\
Privileged basis outlier dimensions & Privileged~\cite{anthropic2023privileged} & \texttt{rounding.rs} & Per-channel scaling \\
Feature steering = causal modification & Scaling~\cite{templeton2024scaling} & SAE steering & Validates gradient-free approach \\
Planning features activate early & Biology~\cite{lindsey2025biology} & \texttt{prefetch.rs} & Predictive RDMA reads \\
Head-type taxonomy (prev-tok, induction) & Framework~\cite{elhage2021framework} & \texttt{rdma\_cache.rs} & Per-head cache policy \\
\bottomrule
\end{tabular}
\end{table}

%=============================================================================
\section{Quantitative Impact Estimates}
\label{sec:impact}
%=============================================================================

\begin{table}[htbp]
\centering
\caption{Estimated bandwidth savings from circuit-aware optimizations on ConnectX-4 100GbE}
\label{tab:impact}
\begin{tabular}{@{}lrrr@{}}
\toprule
\textbf{Optimization} & \textbf{Compression} & \textbf{Effective BW} & \textbf{Stacks With} \\
\midrule
Baseline (dense FP16) & 1$\times$ & 12.5 GB/s & --- \\
MXFP4 quantization (current) & 4$\times$ & 50 GB/s equiv & --- \\
+ Sparse encoding (90\% sparsity) & 10$\times$ & 125 GB/s equiv & MXFP4 \\
+ Delta-mode pipeline & 2--5$\times$ & 250--625 GB/s equiv & Sparse + MXFP4 \\
Combined worst case & 40$\times$ & 500 GB/s equiv & All \\
Combined best case (99\% sparse) & 200$\times$ & 2,500 GB/s equiv & All \\
\bottomrule
\end{tabular}
\end{table}

At the combined best case, the effective bandwidth of a single ConnectX-4 100GbE link approaches that of NVLink, making commodity Ethernet competitive with proprietary interconnects for sparse-activation workloads.

%=============================================================================
\section{Limitations and Caveats}
\label{sec:limitations}
%=============================================================================

\begin{enumerate}
    \item \textbf{SAE reconstruction gap}: The best CLTs achieve only 50\% output-match with the original model and 78.3\% reconstruction accuracy. The remaining computation may resist sparse decomposition.
    \item \textbf{Feature-basis projection adds compute}: Projecting activations into a sparse basis before RDMA transfer requires an encoder forward pass per pipeline stage. Whether the bandwidth savings outweigh the compute cost depends on the network-to-compute ratio.
    \item \textbf{Induction heads are well-studied on toy tasks}: Real-world KV access patterns during long-context generation may be more complex than the induction head model predicts.
    \item \textbf{Composition metrics are static}: The composition strength computed from weights does not account for input-dependent attention pattern variation.
    \item \textbf{No runtime SAE}: Current SAE inference adds latency. The enhancements in Sections~\ref{sec:proposals}.3--\ref{sec:proposals}.5 can be implemented without SAEs using cheaper heuristics (zero-counting for sparsity, weight analysis for head typing).
\end{enumerate}

%=============================================================================
\section{Connection to Thesis Roadmap}
\label{sec:roadmap}
%=============================================================================

The five identified gaps map to existing and new projects in the thesis roadmap:

\begin{table}[htbp]
\centering
\caption{Gap-to-project mapping}
\label{tab:roadmap}
\begin{tabular}{@{}lll@{}}
\toprule
\textbf{Gap} & \textbf{Project} & \textbf{Status} \\
\midrule
1: Composition-aware WFA & P10 (WFA Integration) & Enhancement to existing \\
2: Phase-aware precision & P11 (Precision Pipeline) & Enhancement to existing \\
3: Sparse transfer mode & P12 (Transport Layer) & \textbf{New module needed} \\
4: Head-type prefetching & P13 (Prefetch FSM) & Enhancement to existing \\
5: Delta-mode pipeline & P14 (Pipeline Transport) & \textbf{New module needed} \\
\bottomrule
\end{tabular}
\end{table}

Gaps 1, 2, and 4 require additions to existing \texttt{tensor\_rdma} modules. Gaps 3 and 5 require new functionality in \texttt{transport.rs} (sparse encoding/decoding) and a new delta-compression module.

%=============================================================================
\section{Conclusions}
\label{sec:conclusions}
%=============================================================================

Mechanistic interpretability provides the theoretical foundation for decisions the RDMA middleware currently makes heuristically. The WFA classifier's hot/warm/cold labels correspond to real structural properties: HOT tensors are on composition-critical paths where induction heads demand full KV cache access; COLD tensors carry sparse activations that tolerate aggressive compression.

The five gaps identified---composition awareness, phase-aware precision, sparse transfer encoding, head-type prefetching, and delta-mode pipeline transport---represent concrete, implementable enhancements to the \texttt{tensor\_rdma} crate. The most impactful is sparse transfer mode (Gap~3), which could yield 10--200$\times$ bandwidth savings on activation tensors without any changes to the model itself.

The Transformer Circuits research program is ongoing. As Anthropic extends circuit tracing to larger models and develops faster interpretability tools (MOLT transcoders~\cite{anthropic2025molt} achieve 1000$\times$ fewer FLOPs than standard transcoders), runtime circuit-aware scheduling may become practical for production inference systems.

\bibliographystyle{plain}
\bibliography{update11}

\end{document}
